\documentclass[11pt]{article}

% ─── Packages ─────────────────────────────────────────────────────────────────
\usepackage[margin=1in]{geometry}
\usepackage{amsmath,amssymb,amsthm}
\usepackage{mathtools}
\usepackage{enumitem}
\usepackage{hyperref}
\usepackage{xcolor}
\usepackage{booktabs}

\hypersetup{colorlinks=true, linkcolor=blue!60!black, citecolor=blue!60!black, urlcolor=blue!60!black}

% ─── Theorem environments ─────────────────────────────────────────────────────
\newtheorem{theorem}{Theorem}
\newtheorem{proposition}[theorem]{Proposition}
\newtheorem{lemma}[theorem]{Lemma}
\newtheorem{corollary}[theorem]{Corollary}
\theoremstyle{definition}
\newtheorem{definition}[theorem]{Definition}
\newtheorem{axiom}{Axiom}
\theoremstyle{remark}
\newtheorem{remark}[theorem]{Remark}

% ─── Macros ───────────────────────────────────────────────────────────────────
\newenvironment{ike}{\color{blue}\textbf{Ike:}\\}{}

\newcommand{\KL}{\mathrm{KL}}
\newcommand{\Var}{\mathrm{Var}}
\newcommand{\E}{\mathbb{E}}
\newcommand{\R}{\mathbb{R}}
\DeclareMathOperator{\Cov}{Cov}
\DeclareMathOperator*{\argmax}{arg\,max}
\DeclareMathOperator*{\argmin}{arg\,min}

\title{\textbf{Self-Consistency Constraints on Policy Updates:\\Axiomatic Foundations and Extensions of Mirror Descent for RL}}
\author{Philip Nielsen \and Ike Uchendu \and Kevin He \and Allen Yang\\[4pt]
\textit{CS 2824 --- Spring 2026}}
\date{}

\begin{document}
\maketitle

% ══════════════════════════════════════════════════════════════════════════════
\section{Introduction}
% ══════════════════════════════════════════════════════════════════════════════

Policy optimization in reinforcement learning (RL) faces a fundamental tension: updates should move the policy toward higher expected return, but overly aggressive steps can destabilize learning due to the non-stationarity of the data-generating process. Mirror descent with KL divergence as the Bregman divergence has emerged as one of the most principled frameworks for navigating this tradeoff. It underlies many successful algorithms, including TRPO, MPO, and the recently popularized GRPO.

The standard derivation of the exponential tilt update arrives at it through an optimization lens: the exponential form emerges because KL divergence is the Bregman divergence generated by negative entropy. In this project proposal, we explore an alternative, axiomatic perspective based on a self-consistency property that we call compositionality: updating a trajectory's probability by its total advantage must equal the product of per-step updates. This factorization requirement, combined with a shift-invariance condition, uniquely determines the exponential tilt. The underlying mathematics (Cauchy functional equations for additive-to-multiplicative homomorphisms) is classical and well studied in probability theory and statistical mechanics. The contribution we propose is not the derivation itself, but its consequences when applied to RL policy updates: the constraint that $\eta$ must be a global scalar, a principled variance-based adaptive step-size rule, the characterization of which additional signals (such as advantage uncertainty) can be incorporated self-consistently, and a conjectured connection to entropic risk tilting.

We propose to investigate several extensions of this framework: adaptive step-size selection via a KL budget, extension to discounted rewards, axiomatic constraints on uncertainty-aware tilting signals, and connections to entropic risk measures, Bayesian inference, and multi-agent learning.

% ══════════════════════════════════════════════════════════════════════════════
\section{Background: Mirror Descent for RL}
\label{sec:background}
% ══════════════════════════════════════════════════════════════════════════════

\subsection{The Standard KL-Regularized Derivation}

Consider a Markov decision process with states $s \in \mathcal{S}$, actions $a \in \mathcal{A}$, and a policy $\pi_k$ at iteration $k$. The standard motivation for the KL mirror descent update begins with the trust-region problem:
\begin{equation}\label{eq:trust-region}
  \pi_{k+1} = \argmax_{\pi} \; \E_{s \sim d_{\pi_k}}\!\left[\E_{a \sim \pi(\cdot|s)}[A_k(s,a)]\right]
  \quad \text{s.t.} \quad \E_{s \sim d_{\pi_k}}\!\left[\KL\!\left(\pi(\cdot|s)\,\|\,\pi_k(\cdot|s)\right)\right] \le \Delta,
\end{equation}
where $A_k(s,a) = Q^{\pi_k}(s,a) - V^{\pi_k}(s)$ is the advantage function and $d_{\pi_k}$ is the (discounted) state visitation distribution. Converting the constraint to a Lagrangian penalty with dual variable $\eta_k$, the per-state optimization decomposes into independent problems with closed-form solution:
\begin{equation}\label{eq:mirror-descent}
  \pi_{k+1}(a \mid s) \;\propto\; \pi_k(a \mid s)\,\exp\!\left(\eta_k\, A_k(s,a)\right).
\end{equation}

\begin{ike}
    
\subsubsection*{Deriving the KL Mirror Descent Update}

We want to solve the trust-region problem:
\[
\pi_{k+1} = \arg\max_{\pi} \; \mathbb{E}_{s \sim d_{\pi_k}}\left[\mathbb{E}_{a \sim \pi(\cdot|s)}[A_k(s,a)]\right] \quad \text{s.t.} \quad \mathbb{E}_{s \sim d_{\pi_k}}[\mathrm{KL}(\pi(\cdot|s) \| \pi_k(\cdot|s))] \leq \Delta.
\]

\subsubsection*{Step 1: Form the Lagrangian}

Since we are maximizing over $\pi$, the Lagrangian penalty must \emph{hurt} the objective when the constraint is violated. If $\mathrm{KL}(\pi \| \pi_k) > \Delta$, then $(\mathrm{KL} - \Delta) > 0$, and subtracting a positive quantity from the objective reduces it. So we subtract the penalty term with dual variable $\lambda \geq 0$:
\[
L(\pi, \lambda) = \mathbb{E}_{s \sim d_{\pi_k}}\left[\mathbb{E}_{a \sim \pi(\cdot|s)}[A_k(s,a)]\right] - \lambda\left(\mathbb{E}_{s \sim d_{\pi_k}}[\mathrm{KL}(\pi(\cdot|s) \| \pi_k(\cdot|s))] - \Delta\right).
\]
The dual problem is $\min_{\lambda \geq 0} \max_{\pi} L(\pi, \lambda)$. The outer variable $\lambda$ cranks the penalty up when the constraint is violated; the inner variable $\pi$ maximizes the penalized objective.

\subsubsection*{Step 2: Fix $s$ and decompose}

Since $\pi(\cdot|s)$ at different states don't constrain each other, the inner maximization decomposes into independent per-state problems. For a fixed $s$ and $\lambda$, substituting the definition of KL divergence:
\[
J(\pi) = \sum_a \pi(a|s)\, A_k(s,a) - \lambda \sum_a \pi(a|s) \log \frac{\pi(a|s)}{\pi_k(a|s)}.
\]

\subsubsection*{Step 3: Verify strict concavity}

The first term $\sum_a \pi(a|s) A_k(s,a)$ is a linear function of the variables $\pi(a|s)$, since the advantages $A_k(s,a)$ are fixed coefficients. A linear function has no curvature, so it is both convex and concave.

The KL term is strictly convex in $\pi(\cdot|s)$: it achieves its minimum of zero when $\pi = \pi_k$ and increases as the policies differ, curving upward in every direction. Multiplying by $-\lambda$ (with $\lambda > 0$) flips it to strictly concave.

Therefore $J(\pi)$ is strictly concave (linear + strictly concave = strictly concave), so setting the gradient to zero yields the unique global maximum.

\subsubsection*{Step 4: Take the derivative and solve}

For a specific action $a_i$, take the partial derivative of $J$ with respect to $\pi(a_i|s)$:
\[
\frac{\partial J}{\partial \pi(a_i|s)} = A_k(s,a_i) - \lambda\left[\log \frac{\pi(a_i|s)}{\pi_k(a_i|s)} + 1\right] = 0.
\]
Rearranging:
\[
\log \frac{\pi(a_i|s)}{\pi_k(a_i|s)} = \frac{A_k(s,a_i)}{\lambda} - 1 \quad \implies \quad \pi(a_i|s) = \pi_k(a_i|s) \exp\!\left(\frac{A_k(s,a_i)}{\lambda} - 1\right).
\]
Since this holds for every action $a_i$, we have the full solution for $\pi(\cdot|s)$.

\subsubsection*{Step 5: Normalize}

From Step 4, we have for every action $a_i$:
\[
\pi(a_i|s) = \pi_k(a_i|s) \exp\!\left(\frac{A_k(s,a_i)}{\lambda} - 1\right).
\]
We can separate the exponent:
\[
\pi(a_i|s) = \pi_k(a_i|s) \cdot \exp\!\left(\frac{A_k(s,a_i)}{\lambda}\right) \cdot \exp(-1).
\]
The factor $\exp(-1) = 1/e$ is the same for every action, so it does not affect the relative probabilities between actions. We drop it and normalize $\sum_{a_i} \pi(a_i|s) = 1$ afterward. Writing $\eta_k = 1/\lambda$:
\[
\pi_{k+1}(a|s) \propto \pi_k(a|s) \exp\!\left(\eta_k\, A_k(s,a)\right).
\]

\end{ike}

This is the exponential tilt or Gibbs update, and $\eta_k$ is the inverse temperature. The derivation is clean but fundamentally optimization-centric: the exponential form emerges because KL divergence happens to be the Bregman divergence generated by the negative entropy.

A natural question arises: is there a structural reason why the exponential tilt is the right update, beyond the choice of KL as the divergence?


\begin{ike}
We chose KL divergence to measure distance between policies, and the exponential update fell out of the derivation. But why KL specifically? That choice was not justified since we could have picked a different divergence and gotten a different update form entirely. The question being asked is: is there a deeper reason why the exponential tilt is the right update for policy optimization over probability distributions, independent of the choice of KL?
\end{ike}

\subsection{Prior Work}

Trust-region and mirror-descent methods have a rich history in RL. Natural policy gradient uses the Fisher information metric. TRPO enforces a hard KL constraint; PPO uses a clipped surrogate. MPO and VMPO solve the KL-regularized problem in an E-step/M-step fashion. Recent work on GRPO applies the exponential tilt directly to group-relative advantages in language model fine-tuning.

The core mathematical observation underlying our derivation, that a map commuting with convolution (or equivalently, converting additive signals to multiplicative weights) must be an exponential tilt, is classical. In probability theory, the exponential tilting map is known to be the unique measure-class-preserving map that commutes with convolution; this property is central to large deviations theory. Sandomirskiy and Tamuz (2023) recently gave a rigorous modern treatment, proving that the Boltzmann distribution is the unique family satisfying independence for uncoupled systems. Brandenburger, Bhatt, and Steverson (2019) derived the Boltzmann distribution from different physical axioms using Cauchy functional equations. The Pitman-Koopman-Darmois theorem characterizes exponential families as the unique distributions admitting fixed-dimension sufficient statistics, another manifestation of the same factorization principle.

In the RL context, the most closely related prior work is Path Consistency Learning (PCL) by Nachum et al., which establishes that the optimal policy in entropy-regularized MDPs satisfies a consistency condition along multi-step trajectory segments. Our compositionality axiom is related but distinct: PCL characterizes the fixed point (what the optimal policy looks like), while our axiom constrains the update step (how each iteration must behave). To our knowledge, the specific application of the compositionality axiom to RL policy updates, and the structural consequences it implies (global $\eta$, adaptive step-size rules, admissible uncertainty signals), have not been explicitly developed in the RL literature.

% ══════════════════════════════════════════════════════════════════════════════
\section{Self-Consistency Derivation of Mirror Descent}
\label{sec:self-consistency}
% ══════════════════════════════════════════════════════════════════════════════

We now present the axiomatic derivation. Consider an update rule $w(p, A)$ that assigns an unnormalized weight to an action with prior probability weight $p$ and advantage signal $A \in \R$.

\begin{axiom}[Compositionality]\label{ax:comp}
For any trajectory decomposed into steps with prior weights $p_1, \ldots, p_T$ and per-step advantages $A_1, \ldots, A_T$:
\[
  w\!\left(\prod_{t=1}^T p_t,\; \sum_{t=1}^T A_t\right) = \prod_{t=1}^T w(p_t, A_t).
\]
\end{axiom}

This axiom states that tilting a trajectory by its total advantage is equivalent to independently tilting each step. It is a factorization requirement: the global update decomposes into local updates.

\begin{axiom}[Shift Invariance]\label{ax:neut}
If $A(s,a) = c$ for all actions $a$ at some state $s$ (i.e., the advantage is flat), the normalized policy update produces no change:
\[
  \frac{w(p_a,\, c)}{\sum_{a'} w(p_{a'},\, c)} = p_a \quad \text{for all distributions } (p_a).
\]
\end{axiom}

This axiom states that mirror descent should be insensitive to constant shifts in the tilting signal: only the relative advantage across actions matters.

\begin{theorem}[Exponential Tilt from Compositionality]\label{thm:main}
Under mild measurability assumptions, the unique update rule satisfying Axioms~\ref{ax:comp} and~\ref{ax:neut} is
\[
  w(p, A) = p\,e^{\eta A}
\]
for some constant $\eta \in \R$.
\end{theorem}

\begin{ike}
This theorem says that if your update rule must factorize over time steps (Axiom 1) and must leave the policy unchanged when all actions have equal advantage (Axiom 2), then the \textbf{only} possible update rule is the exponential tilt $w(p, A) = p\,e^{\eta A}$. The exponential form is not an artifact of choosing KL divergence (as we questioned above), it's actually enforced by these two structural requirements.
\end{ike}


\begin{proof}
We proceed in two stages.

\paragraph{Stage 1: Compositionality alone.}
Set all but one factor in Axiom~\ref{ax:comp} to the pair $(1, 0)$. Since $w(1,0)$ must satisfy $w(1,0)^T = w(1,0)$ for all $T$ (by applying the axiom with $T$ identical factors), we have $w(1,0) = 1$. This gives the factorization:

\begin{ike}
We could rewrite this part to specify that the following two results are independent applications of Axiom 1. The factorization $w(p,A) = w(p,0) \cdot w(1,A)$ does not depend on $w(1,0) = 1$.
\end{ike}

\[
  w(p, A) = w(p, 0) \cdot w(1, A).
\]
Define $f(p) \coloneqq w(p, 0)$ and $h(A) \coloneqq w(1, A)$. Substituting back into Axiom~\ref{ax:comp} and separating the $p$ and $A$ dependencies:
\begin{align}
  f(p_1 p_2) &= f(p_1)\,f(p_2), \label{eq:mult-cauchy}\\
  h(A_1 + A_2) &= h(A_1)\,h(A_2). \label{eq:add-cauchy}
\end{align}
Equation~\eqref{eq:mult-cauchy} is the multiplicative Cauchy functional equation; every measurable solution is $f(p) = p^\alpha$ for some $\alpha \in \R$. Equation~\eqref{eq:add-cauchy} is the additive Cauchy exponential equation; every measurable solution is $h(A) = e^{\eta A}$ for some $\eta \in \R$.

Therefore, from compositionality alone:
\begin{equation}\label{eq:general}
  w(p, A) = p^\alpha\, e^{\eta A}.
\end{equation}

\begin{ike}
Question for Philip: The proof claims $w(1,0) = 1$ because $w(1,0) = w(1,0)^T$ for all $T$. But $x = 0$ also satisfies $x = x^T$ for all $T \geq 1$. Is there an assumption that $w$ is not zero?
\end{ike}
Philip: yes, there should be an assumption that $w$ is not zero. Good catch.

\begin{ike}
Define $f(p) := w(p, 0)$ and $h(A) := w(1, A)$. Then:

$f(p_1 p_2) = w(p_1 p_2, 0) = w(p_1, 0) \cdot w(p_2, 0) = f(p_1) \cdot f(p_2)$

$h(A_1 + A_2) = w(1, A_1 + A_2) = w(1, A_1) \cdot w(1, A_2) = h(A_1) \cdot h(A_2)$

And the original factorization becomes:

$w(p, A) = w(p, 0) \cdot w(1, A) = f(p) \cdot h(A)$
\end{ike}

\paragraph{Stage 2: Shift invariance fixes $\alpha = 1$.}

\begin{ike}
Stage 1 gave us $w(p, A) = p^\alpha e^{\eta A}$ with $\alpha$ still free. We need to determine $\alpha$. If $\alpha \neq 1$, the update distorts prior probabilities nonlinearly --- actions get up- or down-weighted based on their probability alone, regardless of advantage. Axiom 2 rules this out by requiring that equal advantages produce no policy change, which forces $\alpha = 1$.
\end{ike}

When $A = c$ for all actions, $w(p, c) = p^\alpha e^{\eta c}$. The normalization constant is $e^{\eta c}\sum_{a'} p_{a'}^\alpha$, so the normalized update is $p^\alpha / \sum_{a'} p_{a'}^\alpha$. Axiom~\ref{ax:neut} requires this to equal $p$ for all distributions. Since $\sum_a p_a = 1$, we need $p^{\alpha - 1} = \sum_{a'} p_{a'}^\alpha$ to be constant in $p$, which forces $\alpha = 1$.
\end{proof}

\begin{ike}
Let's write the normalized update out explicitly so we can see where $\alpha = 1$ comes from.

From Stage 1, the unnormalized weight for a single action $a$ with prior probability $p_a$ and (flat) advantage $c$ is
\[
w(p_a, c) = p_a^\alpha \, e^{\eta c}.
\]
Summing over all actions to get the normalization constant,
\[
Z = \sum_{a'} w(p_{a'}, c) = \sum_{a'} p_{a'}^\alpha \, e^{\eta c} = e^{\eta c} \sum_{a'} p_{a'}^\alpha,
\]
where $e^{\eta c}$ pulls out of the sum because it does not depend on $a'$. The normalized update is therefore
\[
\pi_{\text{new}}(a) = \frac{w(p_a, c)}{Z} = \frac{p_a^\alpha \, e^{\eta c}}{e^{\eta c} \sum_{a'} p_{a'}^\alpha} = \frac{p_a^\alpha}{\sum_{a'} p_{a'}^\alpha}.
\]
The $e^{\eta c}$ factor cancels: when advantages are flat, the tilt multiplies every action by the same constant, so it contributes nothing to the relative weights and disappears under normalization.

Axiom~\ref{ax:neut} now requires $\pi_{\text{new}}(a) = p_a$ for every action $a$ and every starting distribution $(p_a)$:
\[
\frac{p_a^\alpha}{\sum_{a'} p_{a'}^\alpha} = p_a.
\]
Rearranging,
\[
p_a^\alpha = p_a \cdot \sum_{a'} p_{a'}^\alpha \quad\Longrightarrow\quad p_a^{\alpha - 1} = \sum_{a'} p_{a'}^\alpha.
\]
The right-hand side is a single number that does not depend on which action $a$ we picked. The left-hand side depends on $a$ through $p_a$. The only way these can be equal for every action is if $p_a^{\alpha - 1}$ is actually constant in $p_a$, which forces $\alpha - 1 = 0$, i.e.\ $\alpha = 1$.

Sanity check: with $\alpha = 1$, we get $p_a / \sum_{a'} p_{a'} = p_a / 1 = p_a$, since $(p_a)$ is a probability distribution.
\end{ike}

\begin{remark}[Structural consequences]
The derivation reveals several constraints that are not obvious from the optimization perspective:
\begin{enumerate}[label=(\roman*)]
  \item \textbf{Global $\eta$:} The constant $\eta$ is a single scalar shared across all states and time steps. A state-dependent $\eta(s)$ would break the factorization in Axiom~\ref{ax:comp{ike}}.
  \item \textbf{Additive signals enter the exponent:} The derivation assumes the tilting signal is additive (the trajectory signal is $\sum_t A_t$). Any additional per-step statistic that also aggregates additively over the trajectory can be incorporated into the exponent alongside the advantage (see Section~\ref{sec:generalized}).
  \item \textbf{Without neutrality:} The more general family $w(p,A) = p^\alpha e^{\eta A}$ corresponds to $\alpha$-divergence mirror descent (R\'enyi tilting).
\end{enumerate}
\end{remark}

% ══════════════════════════════════════════════════════════════════════════════
\section{Adaptive Step Size}
\label{sec:adaptive-eta}
% ══════════════════════════════════════════════════════════════════════════════

Self-consistency requires $\eta_k$ to be a global scalar, but it can vary across iterations based on global statistics. 

\begin{ike}
Reminder: self-consistency is the name of the framework that is being proposed by this project. How does self-consistency force $\eta_k$ to be a global scalar?

The claim is that if we let $\eta$ depend on the state, Axiom 1 breaks. Let's see this directly.

Suppose $\eta = \eta(s_t)$, so each step gets its own $\eta$ depending on the state visited. Plug $w(p_t, A_t) = p_t \, e^{\eta(s_t) A_t}$ into the right side of Axiom 1:
\[
\prod_{t=1}^{T} w(p_t, A_t) = \left(\prod_{t=1}^{T} p_t\right) \cdot e^{\sum_{t=1}^{T} \eta(s_t) A_t}.
\]
The left side of Axiom 1 is $w\!\left(\prod_t p_t, \sum_t A_t\right)$, which is a function of only two quantities: the product of priors and the sum of advantages. Trajectories that agree on those two quantities must produce the same left side, so they must also produce the same right side.

Now pick two trajectories that both visit states $s_1, s_2$ with priors $p_1, p_2$, but with advantages swapped:
\begin{itemize}
    \item Trajectory X: advantages $A_1, A_2$.
    \item Trajectory Y: advantages $A_2, A_1$.
\end{itemize}
Both have the same $\prod_t p_t = p_1 p_2$ and the same $\sum_t A_t = A_1 + A_2$, so Axiom 1 demands their right-hand sides match:
\[
\eta(s_1) A_1 + \eta(s_2) A_2 = \eta(s_1) A_2 + \eta(s_2) A_1.
\]
Rearranging and factoring:
\[
(\eta(s_1) - \eta(s_2)) A_1 = (\eta(s_1) - \eta(s_2)) A_2.
\]
This has to hold for all advantage values, including $A_1 \neq A_2$. The only way that works is $\eta(s_1) = \eta(s_2)$. Since $s_1$ and $s_2$ were arbitrary, $\eta$ must take the same value at every state --- i.e., it's a global scalar.
\end{ike}

In this section we develop two complementary perspectives on how to set $\eta$: the first (Section~\ref{sec:kl-budget}) derives what $\eta$ should be \emph{proportional to} via a KL budget argument, yielding a variance-normalization rule; the second (Section~\ref{sec:optimal-eta}) derives the \emph{optimal value} of $\eta$ by analyzing the interaction between the mirror descent tilt and downstream distribution shift.

% ──────────────────────────────────────────────────────────────────────────────
\subsection{Adaptive step size via KL Budget}
\label{sec:kl-budget}
% ──────────────────────────────────────────────────────────────────────────────

For the tilted policy $\pi_\eta(a) \propto \pi(a)\,e^{\eta A(a)}$ at a single state, the KL divergence admits the exact expression:
\[
  \KL(\pi_\eta \| \pi) = \eta\,\E_{\pi_\eta}[A] - \log \E_\pi[e^{\eta A}].
\]
Expanding the cumulant generating function $\log \E_\pi[e^{\eta A}]$ to second order:
\begin{equation}\label{eq:kl-expansion}
  \KL(\pi_\eta \| \pi) = \frac{\eta^2}{2}\,\Var_\pi(A) + O(\eta^3).
\end{equation}

\begin{ike}
    Our goal in this section is to choose a step size $\eta$ such that we adhere to some KL budget $\Delta$. That is, we want to solve
    \[
    \Delta = \eta \, \mathbb{E}_{\pi_\eta}[A] - \log \mathbb{E}_\pi[e^{\eta A}]
    \]
    for $\eta$. The problem is that $\eta$ appears on the right-hand side in multiple places (both as the explicit multiplier and inside the tilted policy $\pi_\eta$, and inside the log of the expectation of $e^{\eta A}$), and we can't rearrange this to isolate $\eta$ using algebra. So we need to approximate the right-hand side as a polynomial in $\eta$, which we \emph{can} solve.
    
    The $\log \mathbb{E}_\pi[e^{\eta A}]$ term is the cumulant generating function (CGF) of $A$ evaluated at $\eta$. The CGF has a known Taylor expansion in $\eta$ whose coefficients are the cumulants of $A$ (mean, variance, skewness, ...). Truncating the expansion at second order gives us a polynomial in $\eta$, and the truncation is justified because we expect $\eta$ to be small, so higher-order terms shrink quickly.
\end{ike}

Given a fixed KL budget $\Delta > 0$, we solve $\E_{s \sim d_{\pi_k}}[\KL(\pi_{\eta_k} \| \pi_k)] = \Delta$ using the quadratic approximation:
\begin{equation}\label{eq:eta-rule}
  \boxed{\;\eta_k = \sqrt{\frac{2\,\Delta}{\E_{s \sim d_{\pi_k}}[\Var_{\pi_k}(A_k)]}}\;}
\end{equation}

\begin{ike}
    
\end{ike}

This rule is scale invariant: if rewards are rescaled by $c > 0$, then $A \mapsto cA$, $\Var \mapsto c^2\Var$, and $\eta \mapsto \eta/c$, so the product $\eta A$ is unchanged. The policy update is invariant to the arbitrary scale of the reward function.

Substituting the adaptive rule back into the exponential tilt:
\[
  \pi_{k+1}(a \mid s) \;\propto\; \pi_k(a \mid s)\,\exp\!\left(A_k(s,a)\sqrt{\frac{2\,\Delta}{\E_{s \sim d_{\pi_k}}[\Var_{\pi_k}(A_k)]}}\right)
\]
The KL budget analysis thus does not add a free parameter. In practice, $[\Var_{\pi_k}(A_k)]$ should be estimated using an exponential moving average of $(\hat{A}^\pi(s,a)-\hat{V}^\pi(s))^2$, which does introduce an additional hyperparameter which controls how quickly the normalizer adapts to the environment at the starting of a training run and over the course of the training run. 

\begin{remark}
While empirically, using this tilt is an effective improvement over fixed $\eta$, it still does not capture all of the factors which go into choosing the tilt strength at a particular time step during training. This method does not allow for larger KL step sizes when the tilting signal is clear. Additionally, assuming that $A$ is estimated perfectly, the optimal step size is limited by the state distribution shift, which is not reflected in this update rule. In the next section, we try to capture how $A$ shifts as a function of $\eta$ to determine an optimal value.
\end{remark}

% ──────────────────────────────────────────────────────────────────────────────
\subsection{Adaptive Step Size via One-Step Modeling of Distribution Shift}
\label{sec:optimal-eta}
% ──────────────────────────────────────────────────────────────────────────────

The KL budget analysis treats each state independently. In the MDP, the tilt is applied globally: the policy changes at every state, so the $Q$-values used to compute the tilt become stale. We derive the optimal $\eta$ by expanding the true improvement to second order, accounting for this staleness.

\begin{equation}\label{eq:V_pi_eta}
  V^{\pi_\eta}(s) = V^\pi(s) + \eta\mathrm{Cov}_{a\sim\pi(s)}(Q^{\pi_\eta},A^\pi) + \frac{\eta^2}{2}\mathrm{Cov}_{a\sim\pi(s)}(Q^{\pi_\eta},(A^\pi)^2) + O(\eta^3)
\end{equation}

\begin{align}
  \Delta V(s;\eta) 
  &\coloneqq V^{\pi_\eta}(s) - V^\pi(s) \nonumber\\
  &= \eta\mathrm{Cov}_{a\sim\pi(s)}(Q^{\pi_\eta},A^\pi) + \frac{\eta^2}{2}\mathrm{Cov}_{a\sim\pi(s)}(Q^{\pi_\eta},(A^\pi)^2) + O(\eta^3).
\end{align}
$Q^{\pi_\eta}$ is not known, but we can write it in terms of $Q^\pi$ and $\Delta V(s';\eta)$ to get a recursive equation for $\Delta V(s;\eta)$
\begin{align}
  Q^{\pi_\eta}
  &= \E_{s'|s,a}[r(s,a,s') + \gamma V^{\pi_\eta}(s')]\nonumber\\
  &= \E_{s'|s,a}[r(s,a,s') + \gamma(V^\pi(s') + V^{\pi_\eta}(s') -V^\pi(s'))]\nonumber\\
  &= Q^\pi + \gamma \E_{s'|s,a}[\Delta V(s';\eta)].
\end{align}
Substituting into the covariances:
\begin{align}
    \mathrm{Cov}_{a\sim\pi(s)}(Q^{\pi_\eta},A^\pi) 
    &= \mathrm{Cov}_{a\sim\pi(s)}(Q^\pi,A^\pi) + \gamma\mathrm{Cov}_{a\sim\pi(s)}(\E_{s'|s,a}[\Delta V(s';\eta)],A^\pi)\nonumber\\
    &= \E_{a\sim\pi(s)}[(A^\pi)^2] + \gamma\mathrm{Cov}_{a\sim\pi(s)}(\E_{s'|s,a}[\Delta V(s';\eta)],A^\pi)\nonumber\\
    &= \E_a[(A^\pi)^2] + \eta\gamma\mathrm{Cov}_{a\sim\pi(s)}(\E_{s'|s,a}[\mathrm{Cov}_{a'\sim\pi(s')}(Q^{\pi_\eta},A^\pi)],A^\pi) + O(\eta^2)\nonumber\\
    &= \E_a[(A^\pi)^2] + \eta\gamma\mathrm{Cov}_{a\sim\pi(s)}(\E_{s'|s,a}[\E_{a'}[(A^\pi)^2]],A^\pi) + O(\eta^2)
\end{align}

\begin{align}
    \mathrm{Cov}_{a\sim\pi(s)}(Q^{\pi_\eta},(A^\pi)^2) 
    &= \mathrm{Cov}_{a\sim\pi(s)}(Q^\pi,(A^\pi)^2) + \gamma\mathrm{Cov}_{a\sim\pi(s)}(\E_{s'|s,a}[\Delta V(s';\eta)],(A^\pi)^2)\nonumber\\
    &= \E_{a\sim\pi(s)}[(A^\pi)^3] + O(\eta)
\end{align}
Substituting back into our expression for $\Delta V(s;\eta)$,

\begin{equation}
  \Delta V(s;\eta) = \eta\E_a[(A^\pi)^2] + \eta^2\gamma\mathrm{Cov}_{a\sim\pi(s)}(\E_{s'|s,a}[\E_{a'}[(A^\pi)^2]],A^\pi) + \frac{\eta^2}{2}\E_a[(A^\pi)^3] + O(\eta^3)
\end{equation}
We would like to find the $\eta$ which optimizes this equation, but it cannot depend on $s$ so we take expectations over the state distribution. This gives us three global quantities we need to estimate:
\begin{align}
    v^\pi &\coloneqq \E_{s\sim d_\pi}[\E_a[(A^\pi)^2]]\nonumber\\
    c^\pi &\coloneqq \E_{s\sim d_\pi}[\mathrm{Cov}_{a\sim\pi(s)}(\E_{s'|s,a}[\E_{a'}[(A^\pi)^2]],A^\pi)]\nonumber\\
    \kappa_3^\pi &\coloneqq \E_{s\sim d_\pi}[\E_a[(A^\pi)^3]]\nonumber
\end{align}
This gives
\begin{equation}
  \E_{s\sim d_\pi}[\Delta V(s;\eta)] \approx \eta v^\pi + \eta^2\left(\gamma c^\pi + \frac{\kappa_3^\pi}{2}\right).
\end{equation}
This quantity is optimized at
\begin{equation}\label{eq:eta-star}
  \boxed{\;\eta^* = -\frac{v^\pi}{2\gamma c^\pi + \kappa_3^\pi}\;}
\end{equation}


\paragraph{Practical estimation via dimensionless normalization.}
The three quantities $v^\pi$, $c^\pi$, and $\kappa_3^\pi$ live on different scales: under reward rescaling $r \mapsto \lambda r$, we have
\[
  v^\pi \mapsto \lambda^2 v^\pi,
  \qquad
  c^\pi \mapsto \lambda^3 c^\pi,
  \qquad
  \kappa_3^\pi \mapsto \lambda^3 \kappa_3^\pi.
\]
It is therefore more stable to estimate the \emph{scale} and \emph{shape} of the denominator separately. Define the dimensionless shape statistic
\begin{equation}
  s^\pi
  \;\coloneqq\;
  \frac{2\gamma c^\pi + \kappa_3^\pi}{{v^\pi}^{3/2}}.
\end{equation}
Then the step-size formula can be rewritten as
\begin{equation}
  \eta^*
  \;=\;
  -\frac{v^\pi}{2\gamma c^\pi + \kappa_3^\pi}
  \;=\;
  -\frac{1}{\sqrt{v^\pi}\,s^\pi}.
\end{equation}

In practice, at update step $t$, we form minibatch estimators $\hat v_t$, $\hat c_t$, and $\hat \kappa_{3,t}$ from the current replay batch or on-policy rollout batch, and then construct the dimensionless batch statistic
\begin{equation}
  \hat s_t
  \;\coloneqq\;
  \frac{2\gamma \hat c_t + \hat \kappa_{3,t}}{\hat v_t^{3/2}},
\end{equation}
where
\begin{equation}
  \hat v_t = \frac{1}{B}\sum_i A_i^2,
  \qquad
  \hat \kappa_{3,t} = \frac{1}{B}\sum_i A_i^3,
  \qquad
  \hat c_t = \frac{1}{B}\sum_i (A'_i)^2 \, A_i,
\end{equation}
Where $B$ is the batch size.

Rather than applying an exponential moving average directly to each of the three global quantities, we maintain only two moving averages for the scale and shape:
\begin{align}
  V_t &= \beta_V V_{t-1} + (1-\beta_V)\hat v_t, \\
  S_t &= \beta_S S_{t-1} + (1-\beta_S)\hat s_t.
\end{align}
The running step-size estimate is then reconstructed as
\begin{equation}
  \eta_t
  \;=\;
  -\frac{1}{\sqrt{V_t}\,S_t}.
\end{equation}

This parameterization has two advantages. First, $S_t$ is dimensionless and therefore largely insensitive to the overall reward scale, making it more stationary across training. Second, it allows the variance term $V_t$ and the shape term $S_t$ to be smoothed at different rates: in practice, $V_t$ can usually be tracked with a faster moving average, while $S_t$ should be smoothed more aggressively because it depends on noisier signed third-order statistics. To avoid instability, one should only trust the quadratic step when $S_t$ is safely negative; otherwise it is natural to fall back to the variance-normalized rule from Section~\ref{sec:kl-budget}.

% ──────────────────────────────────────────────────────────────────────────────
\subsection{Adaptive Step Size via Two-Step Modeling of Distribution Shift}
\label{sec:optimal-eta-third-order}
% ──────────────────────────────────────────────────────────────────────────────

The second-order one-step model captures the leading benefit of tilting and the leading penalty from distribution shift. A natural extension is to expand the improvement to third order in $\eta$ and model the induced distribution shift two steps into the future. This gives a cubic surrogate for the expected improvement and therefore a refined adaptive step-size rule.

For any random variable $X(a)$ under the tilted policy
\[
  \pi_\eta(a \mid s)\;\propto\;\pi(a \mid s)\,e^{\eta A^\pi(s,a)},
\]
with $\E_{a\sim\pi(s)}[A^\pi(s,a)] = 0$, the tilted expectation admits the expansion
\begin{align}
  \E_{a\sim\pi_\eta(s)}[X]
  &= \E_{a\sim\pi(s)}[X]
   + \eta\,\Cov_{a\sim\pi(s)}(X,A^\pi) \nonumber\\
  &\quad + \frac{\eta^2}{2}\,\Cov_{a\sim\pi(s)}(X,(A^\pi)^2) \nonumber\\
  &\quad + \frac{\eta^3}{6}\Big(
      \Cov_{a\sim\pi(s)}(X,(A^\pi)^3)
      - 3\,\Cov_{a\sim\pi(s)}(X,A^\pi)\,\E_{a\sim\pi(s)}[(A^\pi)^2]
    \Big)
  + O(\eta^4).
\end{align}
Applying this with $X = Q^{\pi_\eta}(s,a)$ gives
\begin{align}
  \Delta V(s;\eta)
  &\coloneqq V^{\pi_\eta}(s) - V^\pi(s) \nonumber\\
  &= \eta\,T_1(s)
   + \frac{\eta^2}{2}\,T_2(s)
   + \frac{\eta^3}{6}\,T_3(s)
   + O(\eta^4),
\end{align}
where
\begin{align}
  T_1(s)
  &\coloneqq \Cov_{a\sim\pi(s)}(Q^{\pi_\eta},A^\pi), \\
  T_2(s)
  &\coloneqq \Cov_{a\sim\pi(s)}(Q^{\pi_\eta},(A^\pi)^2), \\
  T_3(s)
  &\coloneqq \Cov_{a\sim\pi(s)}(Q^{\pi_\eta},(A^\pi)^3)
  - 3\,\Cov_{a\sim\pi(s)}(Q^{\pi_\eta},A^\pi)\,\E_{a\sim\pi(s)}[(A^\pi)^2].
\end{align}

As before,
\begin{equation}
  Q^{\pi_\eta}(s,a)
  = Q^\pi(s,a) + \gamma\,\E_{s' \mid s,a}[\Delta V(s';\eta)].
\end{equation}
To obtain $\Delta V(s;\eta)$ to order $O(\eta^3)$, we need:
\begin{itemize}
  \item $T_1(s)$ to order $O(\eta^2)$,
  \item $T_2(s)$ to order $O(\eta)$,
  \item $T_3(s)$ to order $O(1)$.
\end{itemize}

Define the local moments
\begin{align}
  m_2(s) &\coloneqq \E_{a\sim\pi(s)}[(A^\pi)^2], \\
  m_3(s) &\coloneqq \E_{a\sim\pi(s)}[(A^\pi)^3], \\
  m_4(s) &\coloneqq \E_{a\sim\pi(s)}[(A^\pi)^4],
\end{align}
and the one-step and two-step shift statistics
\begin{align}
  c_1(s)
  &\coloneqq \Cov_{a\sim\pi(s)}\!\Big(\E_{s' \mid s,a}[m_2(s')],\,A^\pi\Big), \\
  c_2(s)
  &\coloneqq \Cov_{a\sim\pi(s)}\!\Big(\E_{s' \mid s,a}[c_1(s')],\,A^\pi\Big), \\
  d_1(s)
  &\coloneqq \Cov_{a\sim\pi(s)}\!\Big(\E_{s' \mid s,a}[m_3(s')],\,A^\pi\Big), \\
  d_2(s)
  &\coloneqq \Cov_{a\sim\pi(s)}\!\Big(\E_{s' \mid s,a}[m_2(s')],\,(A^\pi)^2\Big).
\end{align}

The second-order one-step approximation at the successor state is
\begin{equation}
  \Delta V(s';\eta)
  =
  \eta\,m_2(s')
  + \eta^2\Big(\gamma\,c_1(s') + \frac{1}{2}m_3(s')\Big)
  + O(\eta^3).
\end{equation}
Substituting this into the three coefficients gives
\begin{align}
  T_1(s)
  &= m_2(s)
   + \eta\,\gamma\,c_1(s)
   + \eta^2\Big(\gamma^2 c_2(s) + \frac{\gamma}{2}d_1(s)\Big)
   + O(\eta^3), \\
  T_2(s)
  &= m_3(s)
   + \eta\,\gamma\,d_2(s)
   + O(\eta^2), \\
  T_3(s)
  &= m_4(s) - 3\,m_2(s)^2 + O(\eta).
\end{align}
Therefore,
\begin{align}
  \Delta V(s;\eta)
  &= \eta\,m_2(s) \nonumber\\
  &\quad + \eta^2\Big(\gamma\,c_1(s) + \frac{1}{2}m_3(s)\Big) \nonumber\\
  &\quad + \eta^3\Big(
      \gamma^2 c_2(s)
      + \frac{\gamma}{2}d_1(s)
      + \frac{\gamma}{2}d_2(s)
      + \frac{1}{6}\big(m_4(s) - 3m_2(s)^2\big)
    \Big)
  + O(\eta^4).
\end{align}

Since $\eta$ must be global, we average over the state distribution and define
\begin{align}
  v^\pi
  &\coloneqq \E_{s\sim d_\pi}[m_2(s)], \\
  c^\pi
  &\coloneqq \E_{s\sim d_\pi}[c_1(s)], \\
  \kappa_3^\pi
  &\coloneqq \E_{s\sim d_\pi}[m_3(s)], \\
  d^\pi
  &\coloneqq \E_{s\sim d_\pi}[c_2(s)], \\
  e^\pi
  &\coloneqq \E_{s\sim d_\pi}[d_1(s)], \\
  f^\pi
  &\coloneqq \E_{s\sim d_\pi}[d_2(s)], \\
  \rho^\pi
  &\coloneqq \E_{s\sim d_\pi}[m_4(s) - 3m_2(s)^2].
\end{align}
Then the expected improvement admits the cubic approximation
\begin{equation}
  \E_{s\sim d_\pi}[\Delta V(s;\eta)]
  \approx
  \eta\,v^\pi
  + \eta^2\,b^\pi
  + \eta^3\,\tau^\pi,
\end{equation}
where
\begin{align}
  b^\pi
  &\coloneqq \gamma\,c^\pi + \frac{1}{2}\kappa_3^\pi, \\
  \tau^\pi
  &\coloneqq \gamma^2 d^\pi
  + \frac{\gamma}{2}e^\pi
  + \frac{\gamma}{2}f^\pi
  + \frac{1}{6}\rho^\pi.
\end{align}

The stationary points satisfy
\begin{equation}
  v^\pi + 2b^\pi \eta + 3\tau^\pi \eta^2 = 0.
\end{equation}
Thus the cubic surrogate yields the step-size
\begin{equation}\label{eq:eta-star-third-order}
  \boxed{
  \eta^*
  =
  \frac{
    -\,b^\pi
    +
    \operatorname{sign}(b^\pi)\sqrt{(b^\pi)^2 - 3\tau^\pi v^\pi}
  }{
    3\tau^\pi
  }}
\end{equation}
where we choose the branch that continuously reduces to the quadratic rule
\[
  \eta^* \to -\frac{v^\pi}{2b^\pi}
  \qquad\text{as}\qquad
  \tau^\pi \to 0.
\]

\paragraph{Practical estimation via scale-shape normalization.}
Under reward rescaling $r \mapsto \lambda r$, we have
\[
  v^\pi \mapsto \lambda^2 v^\pi,
  \qquad
  b^\pi \mapsto \lambda^3 b^\pi,
  \qquad
  \tau^\pi \mapsto \lambda^4 \tau^\pi.
\]
It is therefore natural to separate the overall scale from the two dimensionless shape statistics
\begin{equation}
  s_2^\pi \coloneqq \frac{b^\pi}{(v^\pi)^{3/2}},
  \qquad
  s_3^\pi \coloneqq \frac{\tau^\pi}{(v^\pi)^2}.
\end{equation}
Writing
\[
  x \coloneqq \eta \sqrt{v^\pi},
\]
the cubic surrogate becomes
\begin{equation}
  \E_{s\sim d_\pi}[\Delta V(s;\eta)]
  \approx
  \sqrt{v^\pi}\,\Big(x + s_2^\pi x^2 + s_3^\pi x^3\Big),
\end{equation}
so the optimal dimensionless step $x^*$ solves
\begin{equation}
  1 + 2s_2^\pi x + 3s_3^\pi x^2 = 0.
\end{equation}
Hence
\begin{equation}
  x^*
  =
  \frac{
    -\,s_2^\pi
    +
    \operatorname{sign}(s_2^\pi)\sqrt{(s_2^\pi)^2 - 3s_3^\pi}
  }{
    3s_3^\pi
  },
  \qquad
  \eta^* = \frac{x^*}{\sqrt{v^\pi}}.
\end{equation}

At update step $t$, let $\hat m_{2,i}$, $\hat m_{3,i}$, $\hat m_{4,i}$ denote estimators of
$m_2(s_i)$, $m_3(s_i)$, $m_4(s_i)$, and let
$\hat m'_{2,i}$, $\hat m'_{3,i}$, $\hat c'_{1,i}$ denote estimators of
$m_2(s_i')$, $m_3(s_i')$, $c_1(s_i')$.
Then minibatch estimators of the global coefficients are
\begin{align}
  \hat v_t
  &= \frac{1}{B}\sum_i \hat m_{2,i}, \\
  \hat c_t
  &= \frac{1}{B}\sum_i A_i\,\hat m'_{2,i}, \\
  \hat \kappa_{3,t}
  &= \frac{1}{B}\sum_i \hat m_{3,i}, \\
  \hat d_t
  &= \frac{1}{B}\sum_i A_i\,\hat c'_{1,i}, \\
  \hat e_t
  &= \frac{1}{B}\sum_i A_i\,\hat m'_{3,i}, \\
  \hat f_t
  &= \frac{1}{B}\sum_i \big(A_i^2 - \hat m_{2,i}\big)\,\hat m'_{2,i}, \\
  \hat \rho_t
  &= \frac{1}{B}\sum_i \big(\hat m_{4,i} - 3\hat m_{2,i}^2\big).
\end{align}
We then form
\begin{align}
  \hat b_t
  &= \gamma \hat c_t + \frac{1}{2}\hat \kappa_{3,t}, \\
  \hat \tau_t
  &= \gamma^2 \hat d_t
   + \frac{\gamma}{2}\hat e_t
   + \frac{\gamma}{2}\hat f_t
   + \frac{1}{6}\hat \rho_t,
\end{align}
and the dimensionless batch statistics
\begin{equation}
  \hat s_{2,t} \coloneqq \frac{\hat b_t}{\hat v_t^{3/2}},
  \qquad
  \hat s_{3,t} \coloneqq \frac{\hat \tau_t}{\hat v_t^2}.
\end{equation}

Rather than smoothing all coefficients directly, we maintain moving averages for the scale and shape:
\begin{align}
  V_t &= \beta_V V_{t-1} + (1-\beta_V)\hat v_t, \\
  S_{2,t} &= \beta_2 S_{2,t-1} + (1-\beta_2)\hat s_{2,t}, \\
  S_{3,t} &= \beta_3 S_{3,t-1} + (1-\beta_3)\hat s_{3,t}.
\end{align}
The running step-size estimate is reconstructed by solving
\begin{equation}
  1 + 2S_{2,t}x + 3S_{3,t}x^2 = 0
\end{equation}
for the branch continuous at $S_{3,t}\to 0$:
\begin{equation}
  x_t
  =
  \frac{
    -\,S_{2,t}
    +
    \operatorname{sign}(S_{2,t})\sqrt{S_{2,t}^2 - 3S_{3,t}}
  }{
    3S_{3,t}
  },
  \qquad
  \eta_t = \frac{x_t}{\sqrt{V_t}}.
\end{equation}

In practice, the new third-order terms are materially harder to estimate than the quadratic ones. The quantities $\hat f_t$ and $\hat\rho_t$ require state-conditional second and fourth moments, so a faithful implementation generally needs either multiple action samples per state or auxiliary estimators for $m_2(s)$ and $m_4(s)$. A cheap single-sample proxy can use
\[
  \hat m_{2,i}=A_i^2,\qquad
  \hat m_{3,i}=A_i^3,\qquad
  \hat m_{4,i}=A_i^4,\qquad
  \hat c'_{1,i}=A'_i (A''_i)^2,
\]
but this makes the centered quantities noticeably noisier. As in the second-order case, it is natural to fall back to the quadratic rule or the KL-budget rule whenever the discriminant becomes negative or the resulting step is non-positive.

\subsection*{Second-Order Expansion of $\Delta V$ for Joint $(\eta,\alpha)$}

We now expand the improvement
\[
\Delta V \;:=\; \mathbb{E}_{\pi'}[A] - \mathbb{E}_{\pi}[A]
\]
to second order in $(\eta,\alpha)$, following exactly the same strategy as in the one-step $\eta$ case.

Recall the log-density change
\[
\Delta(a)
=
(\alpha-1)\log \pi(a) + \eta A(a) - \log Z.
\]

Define centered quantities
\[
\tilde A := A - \mathbb{E}[A],
\qquad
\widetilde{\log \pi} := \log \pi - \mathbb{E}[\log \pi],
\]
so that
\[
\Delta \;\approx\; \eta \tilde A + (\alpha-1)\widetilde{\log \pi}.
\]

As before, we use the expansion
\[
\mathbb{E}_{\pi'}[A]
=
\mathbb{E}\big[A e^{\Delta}\big] \big/ \mathbb{E}\big[e^{\Delta}\big],
\]
and expand numerator and denominator to second order.

\subsubsection*{Numerator Expansion}

\[
\mathbb{E}[A e^{\Delta}]
\;\approx\;
\mathbb{E}[A]
+ \mathbb{E}[A \Delta]
+ \frac{1}{2}\mathbb{E}[A \Delta^2].
\]

Compute each term:

\[
\mathbb{E}[A \Delta]
=
\eta\,\mathrm{Var}[A]
+ (\alpha-1)\,\mathrm{Cov}(A,\log \pi),
\]

\[
\mathbb{E}[A \Delta^2]
=
\eta^2 \mathbb{E}[A \tilde A^2]
+ (\alpha-1)^2 \mathbb{E}[A (\widetilde{\log \pi})^2]
+ 2\eta(\alpha-1)\mathbb{E}[A \tilde A\,\widetilde{\log \pi}].
\]

\subsubsection*{Denominator Expansion}

\[
\mathbb{E}[e^{\Delta}]
\;\approx\;
1 + \mathbb{E}[\Delta] + \frac{1}{2}\mathbb{E}[\Delta^2].
\]

Since $\Delta$ is centered to first order,
\[
\mathbb{E}[\Delta] \approx 0,
\]
and
\[
\mathbb{E}[\Delta^2]
=
\eta^2 \mathrm{Var}[A]
+ (\alpha-1)^2 \mathrm{Var}[\log \pi]
+ 2\eta(\alpha-1)\mathrm{Cov}(A,\log \pi).
\]

\subsubsection*{Putting It Together}

Using
\[
\frac{N}{D}
\approx
N \left(1 - \frac{1}{2}\mathbb{E}[\Delta^2]\right),
\]
we obtain
\[
\Delta V
\;\approx\;
\mathbb{E}[A \Delta]
+
\frac{1}{2}\mathbb{E}[A \Delta^2]
-
\frac{1}{2}\mathbb{E}[A]\mathbb{E}[\Delta^2].
\]

\subsubsection*{Final Expression}

Substituting all terms yields
\[
\Delta V
\;\approx\;
\eta\,\mathrm{Var}[A]
+ (\alpha-1)\,\mathrm{Cov}(A,\log \pi)
\]
\[
\quad
+ \frac{1}{2}\eta^2
\Big(
\mathbb{E}[A \tilde A^2]
- \mathbb{E}[A]\mathrm{Var}[A]
\Big)
\]
\[
\quad
+ \frac{1}{2}(\alpha-1)^2
\Big(
\mathbb{E}[A (\widetilde{\log \pi})^2]
- \mathbb{E}[A]\mathrm{Var}[\log \pi]
\Big)
\]
\[
\quad
+ \eta(\alpha-1)
\Big(
\mathbb{E}[A \tilde A\,\widetilde{\log \pi}]
- \mathbb{E}[A]\mathrm{Cov}(A,\log \pi)
\Big).
\]

\subsubsection*{Covariance Form}

This can be written more cleanly in terms of centered moments:
\[
\Delta V
\;\approx\;
\eta\,\mathrm{Var}[A]
+ (\alpha-1)\,\mathrm{Cov}(A,\log \pi)
\]
\[
\quad
+ \frac{1}{2}\eta^2\,\mathrm{Cov}\!\big(A,\tilde A^2\big)
+ \frac{1}{2}(\alpha-1)^2\,\mathrm{Cov}\!\big(A,(\widetilde{\log \pi})^2\big)
\]
\[
\quad
+ \eta(\alpha-1)\,\mathrm{Cov}\!\big(A,\tilde A\,\widetilde{\log \pi}\big).
\]

\subsubsection*{Interpretation}

\begin{itemize}
\item The first-order terms match the one-step result.
\item The second-order terms introduce:
  \begin{itemize}
  \item a curvature correction in the $A$ direction,
  \item a curvature correction in the $\log \pi$ direction,
  \item and a mixed interaction term between the two.
  \end{itemize}
\item Unlike the first-order case, $\alpha$ now contributes non-trivially.
\end{itemize}

\subsection{Adaptive Step Size via Infinite-step Modeling of Distribution Shift}
\label{sec:optimal-eta}

The KL budget analysis treats each state independently. In the full MDP, however, the tilt is applied globally: changing the policy at one state changes which future states are visited, so the $Q$-function used to define the tilt becomes stale. We now derive a second-order approximation to the \emph{true} improvement and optimize that approximation over a single global value of $\eta$.

We begin with the object we want to understand:
\begin{equation}
  \Delta V(s;\eta)
  \;\coloneqq\;
  V^{\pi_\eta}(s) - V^\pi(s),
\end{equation}
where
\[
  \pi_\eta(a\mid s)
  \;\propto\;
  \pi(a\mid s)e^{\eta A^\pi(s,a)}.
\]
By definition,
\begin{align}
  \Delta V(s;\eta)
  &=
  \E_{a\sim \pi_\eta(\cdot\mid s)}\!\bigl[Q^{\pi_\eta}(s,a)\bigr] - V^\pi(s).
\end{align}

We now rewrite $Q^{\pi_\eta}$ in terms of $Q^\pi$ and the downstream value change:
\begin{align}
  Q^{\pi_\eta}(s,a)
  &= \E_{s'|s,a}\!\left[r(s,a,s') + \gamma V^{\pi_\eta}(s')\right] \nonumber\\
  &= \E_{s'|s,a}\!\left[r(s,a,s') + \gamma\bigl(V^\pi(s') + \Delta V(s';\eta)\bigr)\right] \nonumber\\
  &= Q^\pi(s,a) + \gamma \E_{s'|s,a}[\Delta V(s';\eta)].
\end{align}
Define
\begin{equation}
  \phi(s,a;\eta)
  \;\coloneqq\;
  \E_{s'|s,a}[\Delta V(s';\eta)].
\end{equation}
Then
\begin{align}
  \Delta V(s;\eta)
  &=
  \E_{a\sim \pi_\eta(\cdot\mid s)}\!\bigl[Q^\pi(s,a) + \gamma \phi(s,a;\eta)\bigr] - V^\pi(s).
\end{align}

The expectation is under $\pi_\eta$, but $\pi_\eta$ is just the old policy $\pi$ reweighted by $e^{\eta A^\pi}$. Indeed,
\[
  \pi_\eta(a\mid s)
  =
  \frac{\pi(a\mid s)e^{\eta A^\pi(s,a)}}
       {\E_{a'\sim\pi(\cdot\mid s)}[e^{\eta A^\pi(s,a')}]},
\]
so for any function $f(a)$,
\[
  \E_{a\sim\pi_\eta(\cdot\mid s)}[f(a)]
  =
  \frac{
    \E_{a\sim\pi(\cdot\mid s)}[f(a)e^{\eta A^\pi(s,a)}]
  }{
    \E_{a\sim\pi(\cdot\mid s)}[e^{\eta A^\pi(s,a)}]
  }.
\]
Applying this with $f(a)=Q^\pi(s,a)+\gamma\phi(s,a;\eta)$ gives
\begin{align}
  \Delta V(s;\eta)
  &=
  \frac{
    \E_{a\sim\pi(\cdot\mid s)}\!\left[(Q^\pi+\gamma\phi)e^{\eta A^\pi}\right]
  }{
    \E_{a\sim\pi(\cdot\mid s)}[e^{\eta A^\pi}]
  }
  - V^\pi(s).
\end{align}

We now expand the numerator and denominator in powers of $\eta$.

First, since $\E_{a\sim\pi(\cdot\mid s)}[A^\pi(s,a)]=0$,
\begin{align}
  e^{\eta A^\pi}
  &= 1 + \eta A^\pi + \frac{\eta^2}{2}(A^\pi)^2 + O(\eta^3), \\
  \E_\pi[e^{\eta A^\pi}]
  &= 1 + \frac{\eta^2}{2}\E_\pi[(A^\pi)^2] + O(\eta^3).
\end{align}
Here and below, $\E_\pi[\cdot]$ means $\E_{a\sim\pi(\cdot\mid s)}[\cdot]$ with the state $s$ fixed.

Second, because $\Delta V(s;0)=0$, we also have $\phi(s,a;0)=0$. So $\phi$ starts at first order:
\begin{equation}
  \phi(s,a;\eta)
  =
  \eta\,\phi_1(s,a) + \eta^2\,\phi_2(s,a) + O(\eta^3).
\end{equation}
Therefore
\begin{align}
  Q^\pi(s,a) + \gamma\phi(s,a;\eta)
  &=
  Q^\pi(s,a)
  + \eta\gamma\phi_1(s,a)
  + \eta^2\gamma\phi_2(s,a)
  + O(\eta^3).
\end{align}

Now expand the numerator:
\begin{align}
  &(Q^\pi+\gamma\phi)e^{\eta A^\pi} \nonumber\\
  &=
  \bigl(Q^\pi + \eta\gamma\phi_1 + \eta^2\gamma\phi_2\bigr)
  \bigl(1 + \eta A^\pi + \tfrac{\eta^2}{2}(A^\pi)^2\bigr)
  + O(\eta^3) \nonumber\\
  &=
  Q^\pi
  + \eta\bigl(Q^\pi A^\pi + \gamma\phi_1\bigr)
  + \eta^2\bigl(\gamma\phi_2 + \gamma\phi_1A^\pi + \tfrac12 Q^\pi (A^\pi)^2\bigr)
  + O(\eta^3).
\end{align}
Taking expectation under $\pi(\cdot\mid s)$,
\begin{align}
  \E_\pi[(Q^\pi+\gamma\phi)e^{\eta A^\pi}]
  &=
  \E_\pi[Q^\pi] \nonumber\\
  &\quad
  + \eta\Bigl(\E_\pi[Q^\pi A^\pi] + \gamma\E_\pi[\phi_1]\Bigr) \nonumber\\
  &\quad
  + \eta^2\Bigl(
      \gamma\E_\pi[\phi_2]
      + \gamma\E_\pi[\phi_1A^\pi]
      + \tfrac12 \E_\pi[Q^\pi(A^\pi)^2]
    \Bigr)
  + O(\eta^3).
\end{align}

Next, expand the reciprocal of the denominator. Since
\[
  \E_\pi[e^{\eta A^\pi}]
  =
  1 + \frac{\eta^2}{2}\E_\pi[(A^\pi)^2] + O(\eta^3),
\]
we have
\[
  \frac{1}{\E_\pi[e^{\eta A^\pi}]}
  =
  1 - \frac{\eta^2}{2}\E_\pi[(A^\pi)^2] + O(\eta^3).
\]
Multiplying numerator and reciprocal denominator gives
\begin{align}
  \E_{\pi_\eta}[Q^\pi+\gamma\phi]
  &=
  \E_\pi[Q^\pi] \nonumber\\
  &\quad
  + \eta\Bigl(\E_\pi[Q^\pi A^\pi] + \gamma\E_\pi[\phi_1]\Bigr) \nonumber\\
  &\quad
  + \eta^2\Bigl(
      \gamma\E_\pi[\phi_2]
      + \gamma\E_\pi[\phi_1A^\pi]
      + \tfrac12 \E_\pi[Q^\pi(A^\pi)^2]
      - \tfrac12 \E_\pi[Q^\pi]\E_\pi[(A^\pi)^2]
    \Bigr)
  + O(\eta^3).
\end{align}
Subtracting $V^\pi(s)=\E_\pi[Q^\pi]$,
\begin{align}
  \Delta V(s;\eta)
  &=
  \eta\Bigl(\E_\pi[Q^\pi A^\pi] + \gamma\E_\pi[\phi_1]\Bigr) \nonumber\\
  &\quad
  + \eta^2\Bigl(
      \gamma\E_\pi[\phi_2]
      + \gamma\E_\pi[\phi_1A^\pi]
      + \tfrac12 \E_\pi[Q^\pi(A^\pi)^2]
      - \tfrac12 \E_\pi[Q^\pi]\E_\pi[(A^\pi)^2]
    \Bigr)
  + O(\eta^3).
\end{align}

Now simplify each term. Since $Q^\pi = V^\pi + A^\pi$ and $\E_\pi[A^\pi]=0$,
\begin{align}
  \E_\pi[Q^\pi A^\pi]
  &= \E_\pi[(A^\pi)^2], \\
  \E_\pi[\phi_1A^\pi]
  &= \Cov_\pi(\phi_1,A^\pi), \\
  \frac12 \E_\pi[Q^\pi(A^\pi)^2]
  - \frac12 \E_\pi[Q^\pi]\E_\pi[(A^\pi)^2]
  &= \frac12 \Cov_\pi(Q^\pi,(A^\pi)^2) \nonumber\\
  &= \frac12 \E_\pi[(A^\pi)^3].
\end{align}
So
\begin{align}
  \Delta V(s;\eta)
  &=
  \eta\Bigl(\E_\pi[(A^\pi)^2] + \gamma\E_\pi[\phi_1]\Bigr) \nonumber\\
  &\quad
  + \eta^2\Bigl(
      \gamma\E_\pi[\phi_2]
      + \gamma\Cov_\pi(\phi_1,A^\pi)
      + \frac12 \E_\pi[(A^\pi)^3]
    \Bigr)
  + O(\eta^3).
\end{align}

At this point we identify $\phi_1$ and $\phi_2$ from $\Delta V$ itself. Write
\begin{equation}
  \Delta V(s;\eta)
  =
  \eta\,\delta_1(s) + \eta^2\,\delta_2(s) + O(\eta^3).
\end{equation}
Substituting into the definition of $\phi$,
\begin{align}
  \phi(s,a;\eta)
  &= \E_{s'|s,a}[\Delta V(s';\eta)] \nonumber\\
  &= \E_{s'|s,a}\!\left[\eta\,\delta_1(s') + \eta^2\,\delta_2(s') + O(\eta^3)\right] \nonumber\\
  &= \eta\,\E_{s'|s,a}[\delta_1(s')]
   + \eta^2\,\E_{s'|s,a}[\delta_2(s')]
   + O(\eta^3).
\end{align}
Therefore
\begin{align}
  \phi_1(s,a) &= \E_{s'|s,a}[\delta_1(s')], \\
  \phi_2(s,a) &= \E_{s'|s,a}[\delta_2(s')].
\end{align}

Substituting these back in, we get
\begin{align}
  \Delta V(s;\eta)
  &=
  \eta\Bigl(
      \E_{a\sim\pi(s)}[(A^\pi)^2]
      + \gamma \E_{a\sim\pi(s)}\E_{s'|s,a}[\delta_1(s')]
    \Bigr) \nonumber\\
  &\quad
  + \eta^2\Bigl(
      \gamma \E_{a\sim\pi(s)}\E_{s'|s,a}[\delta_2(s')]
      + \gamma \Cov_{a\sim\pi(s)}\!\bigl(\E_{s'|s,a}[\delta_1(s')],A^\pi\bigr)
      + \frac12 \E_{a\sim\pi(s)}[(A^\pi)^3]
    \Bigr)
  + O(\eta^3).
\end{align}
But we also assumed
\[
  \Delta V(s;\eta)
  =
  \eta\,\delta_1(s) + \eta^2\,\delta_2(s) + O(\eta^3),
\]
so we now match coefficients of $\eta$ and $\eta^2$.

Matching the $O(\eta)$ terms gives
\begin{equation}\label{eq:delta1-correct}
  \delta_1(s)
  =
  \E_{a\sim\pi(s)}[(A^\pi)^2]
  + \gamma \E_{a\sim\pi(s)}\E_{s'|s,a}[\delta_1(s')].
\end{equation}
Matching the $O(\eta^2)$ terms gives
\begin{equation}\label{eq:delta2-correct}
  \delta_2(s)
  =
  \frac12 \E_{a\sim\pi(s)}[(A^\pi)^3]
  + \gamma \Cov_{a\sim\pi(s)}\!\bigl(\E_{s'|s,a}[\delta_1(s')],A^\pi\bigr)
  + \gamma \E_{a\sim\pi(s)}\E_{s'|s,a}[\delta_2(s')].
\end{equation}

These are Bellman-style equations for the first- and second-order coefficients of the improvement. To make the structure more compact, define the one-step source terms
\begin{align}
  u_1(s)
  &\coloneqq
  \E_{a\sim\pi(s)}[(A^\pi)^2]
  =
  \Var_{a\sim\pi(s)}(A^\pi), \\
  u_2(s)
  &\coloneqq
  \frac12 \E_{a\sim\pi(s)}[(A^\pi)^3]
  + \gamma \Cov_{a\sim\pi(s)}\!\bigl(\E_{s'|s,a}[\delta_1(s')],A^\pi\bigr).
\end{align}
Also define the one-step state-transition operator under $\pi$ by
\begin{equation}
  (P^\pi f)(s)
  \;\coloneqq\;
  \E_{a\sim\pi(s)}\E_{s'|s,a}[f(s')].
\end{equation}
Then the two equations become
\begin{align}
  \delta_1 &= u_1 + \gamma P^\pi \delta_1, \\
  \delta_2 &= u_2 + \gamma P^\pi \delta_2.
\end{align}
So
\begin{align}
  \delta_1 &= (I-\gamma P^\pi)^{-1}u_1, \\
  \delta_2 &= (I-\gamma P^\pi)^{-1}u_2.
\end{align}

We now choose a \emph{single global} value of $\eta$. To do that, fix a start-state distribution $\rho$ and optimize the average second-order surrogate
\begin{align}
  J_\rho(\eta)
  &\coloneqq
  \E_{s_0\sim\rho}[\Delta V(s_0;\eta)] \nonumber\\
  &=
  \eta\,\E_\rho[\delta_1]
  + \eta^2\,\E_\rho[\delta_2]
  + O(\eta^3).
\end{align}
Define the discounted state visitation distribution induced by $\rho$ and $\pi$:
\begin{equation}
  d_\rho^\pi
  \;\coloneqq\;
  (1-\gamma)\sum_{t=0}^\infty \gamma^t\,\rho(P^\pi)^t.
\end{equation}
Then for any bounded function $u$,
\begin{align}
  \E_\rho[(I-\gamma P^\pi)^{-1}u]
  &= \sum_{t=0}^\infty \gamma^t \E_\rho[(P^\pi)^t u] \nonumber\\
  &= \sum_{t=0}^\infty \gamma^t \E_{s_t\sim \rho(P^\pi)^t}[u(s_t)] \nonumber\\
  &= \frac{1}{1-\gamma}\E_{s\sim d_\rho^\pi}[u(s)].
\end{align}
Applying this with $u=u_1$ and $u=u_2$ yields
\begin{align}
  \E_\rho[\delta_1]
  &= \frac{1}{1-\gamma}\E_{s\sim d_\rho^\pi}[u_1(s)], \\
  \E_\rho[\delta_2]
  &= \frac{1}{1-\gamma}\E_{s\sim d_\rho^\pi}[u_2(s)].
\end{align}
So
\begin{align}
  J_\rho(\eta)
  &=
  \frac{1}{1-\gamma}
  \Bigl(
    \eta\,\bar u_1 + \eta^2\,\bar u_2
  \Bigr)
  + O(\eta^3),
\end{align}
where
\begin{align}
  \bar u_1 &\coloneqq \E_{s\sim d_\rho^\pi}[u_1(s)], \\
  \bar u_2 &\coloneqq \E_{s\sim d_\rho^\pi}[u_2(s)].
\end{align}
Writing these out explicitly,
\begin{align}
  \bar u_1
  &=
  \E_{s\sim d_\rho^\pi}\!\left[\Var_{a\sim\pi(s)}(A^\pi)\right], \\
  \bar u_2
  &=
  \E_{s\sim d_\rho^\pi}\!\left[
    \frac12 \E_{a\sim\pi(s)}[(A^\pi)^3]
    + \gamma \Cov_{a\sim\pi(s)}\!\bigl(\E_{s'|s,a}[\delta_1(s')],A^\pi\bigr)
  \right].
\end{align}

Ignoring the overall factor $(1-\gamma)^{-1}$, the quadratic surrogate is
\begin{equation}
  J_{\mathrm{quad}}(\eta)
  \;\coloneqq\;
  \eta\,\bar u_1 + \eta^2\,\bar u_2.
\end{equation}
Differentiating,
\[
  J_{\mathrm{quad}}'(\eta)
  =
  \bar u_1 + 2\eta\bar u_2.
\]
So the maximizer of the quadratic approximation is
\begin{equation}\label{eq:eta-star-correct}
  \boxed{
  \eta^*_{\mathrm{quad}}
  =
  -\frac{\bar u_1}{2\bar u_2}
  }.
\end{equation}
Substituting the definitions of $\bar u_1$ and $\bar u_2$ gives
\begin{equation}
  \boxed{
  \eta^*_{\mathrm{quad}}
  =
  -\frac{
    \E_{s\sim d_\rho^\pi}\!\left[\Var_{a\sim\pi(s)}(A^\pi)\right]
  }{
    \E_{s\sim d_\rho^\pi}\!\left[\E_{a\sim\pi(s)}[(A^\pi)^3]\right]
    + 2\gamma\,
    \E_{s\sim d_\rho^\pi}\!\left[
      \Cov_{a\sim\pi(s)}\!\bigl(\E_{s'|s,a}[\delta_1(s')],A^\pi\bigr)
    \right]
  }.
  }
\end{equation}

Thus, the first-order gain from increasing $\eta$ is controlled by the advantage variance, while the second-order correction has two parts: a skewness term and a distribution-shift term. The distribution-shift term captures the fact that actions that look good under the old critic can move the agent into states where that same critic is already becoming stale, which penalizes overly aggressive tilts.

% ══════════════════════════════════════════════════════════════════════════════
\section{Proposed Extensions}
\label{sec:extensions}
% ══════════════════════════════════════════════════════════════════════════════

\subsection{Generalized Self-Consistency and Admissible Tilting Signals}
\label{sec:generalized}

The compositionality axiom as stated in Section~\ref{sec:self-consistency} uses the advantage as the sole tilting signal. But in practice, an agent may have access to additional per-step statistics, such as an estimate of the advantage variance $\sigma^2(s,a)$. The self-consistency framework naturally accommodates multiple signals, provided each aggregates additively over the trajectory.

\paragraph{Joint update with advantage and uncertainty.}
Suppose the update depends on both an advantage signal $A$ (aggregating additively) and a variance signal $\sigma^2$ (also aggregating additively, since $\Var(\sum_t A_t) = \sum_t \Var(A_t)$ under conditional independence). Self-consistency with both signals simultaneously gives $w(p, A, \sigma^2) = f(p) \cdot h_1(A) \cdot h_2(\sigma^2)$. The three functional equations and their solutions are:
\begin{align*}
  f(p_1 p_2) &= f(p_1)\,f(p_2) &&\Longrightarrow\quad f(p) = p^\alpha, \\
  h_1(A_1 + A_2) &= h_1(A_1)\,h_1(A_2) &&\Longrightarrow\quad h_1(A) = e^{c_1 A}, \\
  h_2(\sigma_1^2 + \sigma_2^2) &= h_2(\sigma_1^2)\,h_2(\sigma_2^2) &&\Longrightarrow\quad h_2(\sigma^2) = e^{c_2\,\sigma^2}.
\end{align*}
Applying shift invariance ($\alpha = 1$), the unique self-consistent joint update is:
\begin{equation}\label{eq:joint-update}
  w(p, A, \sigma^2) = p \cdot \exp\!\left(c_1\, A + c_2\,\sigma^2\right).
\end{equation}
At first glance this resembles a UCB/LCB bonus, but with variance rather than standard deviation, and the two terms have conflicting units (reward vs.\ reward$^2$). This is not a problem because everything inside an exponent must be dimensionless: $c_1$ has units $[\text{reward}]^{-1}$ and $c_2$ has units $[\text{reward}]^{-2}$. But it does raise the question: what determines the relationship between $c_1$ and $c_2$? We revisit this in Research Direction~(2) below.

\paragraph{Implications.}
The self-consistency framework constrains not only the functional form of the update but also which additional signals are admissible: any per-step statistic that aggregates additively over the trajectory can enter the exponent with its own coefficient. Variance is a natural candidate because it is additive under conditional independence. The framework thus provides a principled way to incorporate uncertainty into the policy update, rather than relying on ad hoc bonus terms.

\subsection{Additional Research Directions}

\begin{enumerate}[label=\textbf{(\arabic*)}]

\item \textbf{R\'enyi Tilting and the $\alpha$-Family.}
Without the neutrality axiom, the general solution is $w(p,A) = p^\alpha e^{\eta A}$, which corresponds to $\alpha$-divergence mirror descent. The case $\alpha \neq 1$ yields R\'enyi-divergence trust regions. We propose to characterize the full $(\alpha, \eta)$ family: what optimization problem does $\alpha \neq 1$ solve? How does it relate to risk-sensitive RL objectives? Can the $\alpha$ parameter serve as a knob for controlling exploration vs.\ exploitation in a principled way?

\item \textbf{Cumulants, Entropic Risk, and Uncertain Advantages.}
The joint update in Eq.~\eqref{eq:joint-update} involves the first two cumulants ($\kappa_1 = A$, $\kappa_2 = \sigma^2$) of the advantage belief, each entering with a free coefficient. But all cumulants are additive over independent random variables (this is their defining property), so the self-consistency framework extends to arbitrarily many cumulants, yielding $w = p \cdot \exp(\sum_n c_n \kappa_n)$ with free $c_n$. A natural question is whether an additional coherence requirement can fix the $c_n$. We conjecture that requiring the multi-cumulant update to equal the expected exponential tilt under the advantage belief, $\E_A[pe^{\eta A}] = p\,\exp(K_A(\eta))$ where $K_A$ is the cumulant generating function, forces $c_n = \eta^n / n!$ and recovers entropic risk tilting $\pi'(a|s) \propto \pi(a|s)\,\E[e^{\eta A(s,a)}]$. If correct, this would resolve the units issue (each $\eta^n \kappa_n / n!$ is dimensionless), collapse the infinite family of free coefficients to a single parameter $\eta$, and provide a self-consistency derivation of entropic risk measures in RL. We propose to formalize and prove this connection.

\item \textbf{Extension to Discounted Rewards.}
In the discounted setting, the trajectory-level signal is $G(\tau) = \sum_t \gamma^t A(s_t, a_t)$, which is additive but with per-step contributions $\tilde{A}_t = \gamma^t A_t$. Self-consistency then yields per-step updates with a time-dependent effective step size $\eta_t = \eta\,\gamma^t$. This structured time-dependence is forced by the discount factor, not a free design choice, and implies that earlier steps receive more aggressive updates. We propose to investigate whether this time-dependent $\eta_t$ can be reconciled with stationary policy representations, or whether it motivates time-indexed or horizon-aware policy architectures.

\item \textbf{Empirical Validation of the Adaptive $\eta$ Rule.}
The variance-based adaptive $\eta$ rule in Eq.~\eqref{eq:eta-rule} has not, to our knowledge, been tested empirically in deep RL. We propose to implement it in a continuous-control setting (e.g., MuJoCo) and compare against fixed $\eta$, linearly decaying $\eta$, and the PPO clipping mechanism. Key metrics: final performance, stability across seeds, and sensitivity to reward scale.

\item \textbf{Multi-Agent and Game-Theoretic Extensions.}
In multi-agent settings, the ``trajectory'' factorizes across both time steps and agents. The compositionality axiom can be extended to require factorization across agents as well as time. This would constrain the joint update rule and may shed light on when independent learning is self-consistent and when it is not. We propose to derive the axiomatic constraints for the multi-agent case and connect them to existing results on independent vs.\ centralized learning.

\item \textbf{Connection to Bayesian Inference and Posterior Updates.}
The exponential tilt $\pi' \propto \pi\,e^{\eta A}$ has the same form as a Bayesian posterior update with log-likelihood $\eta A$. We propose to make this connection precise: under what model does the advantage function play the role of a log-likelihood? This may connect to the ``RL as inference'' framework and provide a new justification for why self-consistent updates are well-calibrated in a Bayesian sense.

\end{enumerate}

% ══════════════════════════════════════════════════════════════════════════════


\end{document}
